% answers.tex —— 1985 年·全国卷·文科 · 参考答案与详细解析片段
% 由 tools/generate_projects.py 依来源总账自 1985/1985_数学真题_全国文_试题.md + 1985/1985_数学真题_全国文_解析正文.md + 1985/1985_数学真题_全国文_答案要点.md 生成，勿手改（改源 MD 后重跑）


\begin{gktip}
固定源只有一页题面，不含答案。以下逐题给出本库独立推导；答案结论不宣称为官方发布版本。
\end{gktip}

\begin{answer}
\gksolhead{1}{}

\textbf{题目}：

\Needspace{5\baselineskip}
如果正方体 $ABCD-A'B'C'D'$ 的棱长为 $a$，那么四面体 $A'-ABD$ 的体积是\mbox{（　　）}

\nopagebreak[4]
\[
\text{（A）}\ \frac{a^3}{2}
\qquad
\text{（B）}\ \frac{a^3}{3}
\qquad
\text{（C）}\ \frac{a^3}{4}
\qquad
\text{（D）}\ \frac{a^3}{6}
\]

\Needspace{4\baselineskip}
\textbf{解析}

\Needspace{5\baselineskip}
三角形 $ABD$ 是边长为 $a$ 的等边三角形，面积为 $\frac12a^2$；点 $A'$ 到底面 $ABD$ 的高为 $AA'=a$。所以

\nopagebreak[4]
\[
V_{A'-ABD}=\frac13\cdot\frac12a^2\cdot a=\frac{a^3}{6}.
\]

故选（D）。

\gkanswer{答案：（D）$\dfrac{a^3}{6}$}

\end{answer}
\begin{answer}
\gksolhead{2}{}

\textbf{题目}：

$\tan x=1$ 是 $x=\dfrac54\pi$ 的\mbox{（　　）}

（A）必要条件　　（B）充分条件　　（C）充分必要条件　　（D）既不充分又不必要的条件

\Needspace{4\baselineskip}
\textbf{解析}

\Needspace{5\baselineskip}
$x=\frac54\pi$ 能推出 $\tan x=1$，但 $\tan x=1$ 的解为

\nopagebreak[4]
\[
x=\frac\pi4+k\pi,\qquad k\in\mathbb Z,
\]

并不必然等于 $\frac54\pi$。因此“$\tan x=1$”是“$x=\frac54\pi$”的必要条件，选（A）。

\gkanswer{答案：（A）必要条件}

\end{answer}
\begin{answer}
\gksolhead{3}{}

\textbf{题目}：

\Needspace{5\baselineskip}
设集合

\nopagebreak[4]
\[
X=\{0,1,2,4,5,7\},\qquad
Y=\{1,3,6,8,9\},\qquad
Z=\{3,7,8\},
\]

\Needspace{5\baselineskip}
那么集合 $(X\cap Y)\cup Z$ 是\mbox{（　　）}

\nopagebreak[4]
\[
\text{（A）}\ \{0,1,2,6,8\}
\qquad
\text{（B）}\ \{3,7,8\}
\]

\nopagebreak[4]
\[
\text{（C）}\ \{1,3,7,8\}
\qquad
\text{（D）}\ \{1,3,6,7,8\}.
\]

\Needspace{4\baselineskip}
\textbf{解析}

\nopagebreak[4]
\[
X\cap Y=\{1\},
\]

\Needspace{5\baselineskip}
所以

\nopagebreak[4]
\[
(X\cap Y)\cup Z=\{1\}\cup\{3,7,8\}=\{1,3,7,8\}.
\]

故选（C）。

\gkanswer{答案：（C）$\{1,3,7,8\}$}

\end{answer}
\begin{answer}
\gksolhead{4}{}

\textbf{题目}：

在下面给出的函数中，哪一个函数既是区间 $\left(0,\dfrac\pi2\right)$ 上的增函数又是以 $\pi$ 为周期的偶函数\mbox{（　　）}

（A）$y=x^2\ (x\in\mathbb R)$

（B）$y=|\sin x|\ (x\in\mathbb R)$

（C）$y=\cos2x\ (x\in\mathbb R)$

（D）$y=e^{\sin2x}\ (x\in\mathbb R)$

\Needspace{4\baselineskip}
\textbf{解析}

\Needspace{5\baselineskip}
$|\sin x|$ 是偶函数，且

\nopagebreak[4]
\[
|\sin(x+\pi)|=|-\sin x|=|\sin x|,
\]

所以它以 $\pi$ 为周期。在 $\left(0,\frac\pi2\right)$ 上，$\sin x$ 严格递增，故 $|\sin x|=\sin x$ 也严格递增。

其余选项分别因不具 $\pi$ 周期、在该区间递减或不为偶函数而不合题意。故选（B）。

\gkanswer{答案：（B）$y=|\sin x|$}

\end{answer}
\begin{answer}
\gksolhead{5}{}

\textbf{题目}：

用 $1,2,3,4,5$ 这五个数字，可以组成比 $20000$ 大，并且百位数不是数字 $3$ 的没有重复数字的五位数，共有\mbox{（　　）}

（A）96 个　　（B）78 个　　（C）72 个　　（D）64 个

\Needspace{4\baselineskip}
\textbf{解析}

五位数必须由 $1,2,3,4,5$ 各使用一次。首位不能为 $1$，且要大于 $20000$，故首位可为 $2,3,4,5$。

\Needspace{5\baselineskip}
首位为 $2\text{、}4\text{、}5$ 时，剩余数字中有 $3$，百位不能取 $3$，百位有 $3$ 种取法，其余两位和千位排列有 $3!$ 种，共

\nopagebreak[4]
\[
3\times3\times3!=54
\]

种。

\Needspace{5\baselineskip}
首位为 $3$ 时，剩余数字中没有 $3$，百位有 $4$ 种取法，其余三位排列有 $3!$ 种，共

\nopagebreak[4]
\[
4\times3!=24
\]

\Needspace{5\baselineskip}
种。因此总数为

\nopagebreak[4]
\[
54+24=78.
\]

故选（B）。

\gkanswer{答案：（B）$78$ 个}

\end{answer}
\begin{answer}
\gksolhead{6}{}

\textbf{题目}：

\Needspace{5\baselineskip}
求函数

\nopagebreak[4]
\[
y=\frac{\sqrt{4-x^2}}{x-1}
\]

的定义域。

\Needspace{4\baselineskip}
\textbf{解析}

\Needspace{5\baselineskip}
根式要求 $4-x^2\ge0$，即 $-2\le x\le2$；分母要求 $x\ne1$。因此定义域为

\nopagebreak[4]
\[
[-2,1)\cup(1,2].
\]

\gkanswer{答案：$[-2,1)\cup(1,2]$}

\end{answer}
\begin{answer}
\gksolhead{7}{}

\textbf{题目}：

\Needspace{5\baselineskip}
求圆锥曲线

\nopagebreak[4]
\[
3x^2-y^2+6x+2y-1=0
\]

的离心率。

\Needspace{4\baselineskip}
\textbf{解析}

\Needspace{5\baselineskip}
配方得

\nopagebreak[4]
\[
3(x^2+2x)-(y^2-2y)-1=0,
\]

\Needspace{5\baselineskip}
即

\nopagebreak[4]
\[
3(x+1)^2-(y-1)^2=3,
\]

\Needspace{5\baselineskip}
或

\nopagebreak[4]
\[
\frac{(x+1)^2}{1^2}-\frac{(y-1)^2}{3}=1.
\]

\Needspace{5\baselineskip}
这是焦点在 $x$ 轴方向的双曲线，$a^2=1$、$b^2=3$，故

\nopagebreak[4]
\[
c^2=a^2+b^2=4,\qquad e=\frac ca=2.
\]

\gkanswer{答案：$e=2$}

\end{answer}
\begin{answer}
\gksolhead{8}{}

\textbf{题目}：

\Needspace{5\baselineskip}
求函数

\nopagebreak[4]
\[
y=-x^2+4x-2
\]

在区间 $[0,3]$ 上的最大值和最小值。

\Needspace{4\baselineskip}
\textbf{解析}

\nopagebreak[4]
\[
y=-x^2+4x-2=-(x-2)^2+2.
\]

\Needspace{5\baselineskip}
抛物线顶点横坐标 $x=2$ 在区间 $[0,3]$ 内，因此最大值为 $2$。端点值为

\nopagebreak[4]
\[
y(0)=-2,\qquad y(3)=1,
\]

故最小值为 $-2$。

\gkanswer{答案：最大值为 $2$，最小值为 $-2$}

\end{answer}
\begin{answer}
\gksolhead{9}{}

\textbf{题目}：

\Needspace{5\baselineskip}
设

\nopagebreak[4]
\[
(3x-1)^6=a_6x^6+a_5x^5+a_4x^4+a_3x^3+a_2x^2+a_1x+a_0,
\]

求 $a_6+a_5+a_4+a_3+a_2+a_1+a_0$ 的值。

\Needspace{4\baselineskip}
\textbf{解析}

\Needspace{5\baselineskip}
令 $x=1$，则

\nopagebreak[4]
\[
a_6+a_5+a_4+a_3+a_2+a_1+a_0=(3-1)^6=2^6=64.
\]

\gkanswer{答案：$64$}

\end{answer}
\begin{answer}
\gksolhead{10}{}

\textbf{题目}：

设 $i$ 是虚数单位，求 $(1+i)^6$ 的值。

\Needspace{4\baselineskip}
\textbf{解析}

\nopagebreak[4]
\[
(1+i)^2=2i,\qquad
(1+i)^4=-4,
\]

\Needspace{5\baselineskip}
所以

\nopagebreak[4]
\[
(1+i)^6=(1+i)^4(1+i)^2=-4\cdot2i=-8i.
\]

\gkanswer{答案：$-8i$}

\end{answer}
\begin{answer}
\gksolhead{11}{}

\textbf{题目}：

\Needspace{5\baselineskip}
设

\nopagebreak[4]
\[
S_1=1^2,\qquad
S_2=1^2+2^2+1^2,\qquad
S_3=1^2+2^2+3^2+2^2+1^2,\quad\ldots,
\]

\nopagebreak[4]
\[
S_n=1^2+2^2+3^2+\cdots+n^2+\cdots+3^2+2^2+1^2.
\]

\Needspace{5\baselineskip}
用数学归纳法证明：公式

\nopagebreak[4]
\[
S_n=\frac{n(2n^2+1)}{3}
\]

对所有的正整数 $n$ 都成立。

\Needspace{4\baselineskip}
\textbf{解析}

\Needspace{5\baselineskip}
当 $n=1$ 时，

\nopagebreak[4]
\[
S_1=1=\frac{1(2\cdot1^2+1)}3.
\]

\Needspace{5\baselineskip}
假设当 $n=k$ 时成立，即

\nopagebreak[4]
\[
S_k=\frac{k(2k^2+1)}3.
\]

\Needspace{5\baselineskip}
由 $S_{k+1}$ 比 $S_k$ 两侧各增加一项 $(k+1)^2$，有

\nopagebreak[4]
\[
\begin{aligned}
S_{k+1}
&=S_k+2(k+1)^2\\
&=\frac{k(2k^2+1)}3+2(k+1)^2\\
&=\frac{2k^3+6k^2+7k+6}{3}\\
&=\frac{(k+1)\bigl(2(k+1)^2+1\bigr)}3.
\end{aligned}
\]

\Needspace{5\baselineskip}
因此由数学归纳法，公式

\nopagebreak[4]
\[
S_n=\frac{n(2n^2+1)}3
\]

对所有正整数 $n$ 成立。

\gkanswer{答案：$S_n=\dfrac{n(2n^2+1)}{3}$}

\end{answer}
\begin{answer}
\gksolhead{12}{}

\textbf{题目}：

\Needspace{5\baselineskip}
证明三角恒等式：

\nopagebreak[4]
\[
2\sin^4x+\frac34\sin^22x+5\cos^4x-\cos3x\cos x
=2(1+\cos^2x).
\]

\Needspace{4\baselineskip}
\textbf{解析}

利用 $\sin^22x=4\sin^2x\cos^2x$ 和
\Needspace{5\baselineskip}
$\cos3x\cos x=(4\cos^3x-3\cos x)\cos x$，左端为

\nopagebreak[4]
\[
2\sin^4x+3\sin^2x\cos^2x+5\cos^4x-4\cos^4x+3\cos^2x.
\]

\Needspace{5\baselineskip}
令 $s=\sin^2x$、$c=\cos^2x$，则 $s+c=1$，上式等于

\nopagebreak[4]
\[
2s^2+3sc+c^2+3c
=2(1-c)^2+3(1-c)c+c^2+3c
=2+2c
=2(1+\cos^2x).
\]

恒等式得证。

\gkanswer{答案：恒等式成立}

\end{answer}
\begin{answer}
\gksolhead{13}{}

\textbf{题目}：

\Needspace{5\baselineskip}
解方程：

\nopagebreak[4]
\[
\lg(3-x)-\lg(3+x)=\lg(1-x)-\lg(2x+1).
\]

\Needspace{4\baselineskip}
\textbf{解析}

\Needspace{5\baselineskip}
定义域要求

\nopagebreak[4]
\[
3-x>0,\quad 3+x>0,\quad 1-x>0,\quad 2x+1>0,
\]

\Needspace{5\baselineskip}
所以 $-1<x<1$。在此范围内合并对数得

\nopagebreak[4]
\[
\frac{3-x}{3+x}=\frac{1-x}{2x+1}.
\]

\Needspace{5\baselineskip}
交叉相乘：

\nopagebreak[4]
\[
(3-x)(2x+1)=(3+x)(1-x),
\]

\Needspace{5\baselineskip}
化简为

\nopagebreak[4]
\[
x(x-7)=0.
\]

只有 $x=0$ 在定义域内，故方程的解为 $x=0$。

\gkanswer{答案：$x=0$}

\end{answer}
\begin{answer}
\gksolhead{14}{}

\textbf{题目}：

\Needspace{5\baselineskip}
解不等式：

\nopagebreak[4]
\[
\sqrt{2x+5}>x+1.
\]

\Needspace{4\baselineskip}
\textbf{解析}

定义域为 $x\ge-\frac52$。当 $x<-1$ 时，右端 $x+1<0$，而左端非负，故不等式自动成立。

\Needspace{5\baselineskip}
当 $x\ge-1$ 时两边非负，可以平方：

\nopagebreak[4]
\[
2x+5>(x+1)^2
\quad\Longleftrightarrow\quad
x^2<4.
\]

\Needspace{5\baselineskip}
与 $x\ge-1$ 合并得 $-1\le x<2$。因此总解集为

\nopagebreak[4]
\[
\left[-\frac52,-1\right)\cup[-1,2)
=\left[-\frac52,2\right).
\]

\gkanswer{答案：$x\in\left[-\dfrac52,2\right)$}

\end{answer}
\begin{answer}
\gksolhead{15}{}

\textbf{题目}：

已知三棱锥 $V-ABC$ 的三个侧面与底面所成的二面角都是 $\beta$，它的高是 $h$。求这个棱锥底面的内切圆半径。

\Needspace{4\baselineskip}
\textbf{解析}

\Needspace{5\baselineskip}
设高的垂足为 $O$。三个侧面与底面的二面角相等，说明从 $O$ 到底面三边的距离相等，因此 $O$ 是底面三角形的内心。设内切圆半径为 $r$，在高 $VO$ 与任一底边的垂直截面内，得到直角三角形，其中

\nopagebreak[4]
\[
\tan\beta=\frac hr.
\]

\Needspace{5\baselineskip}
所以

\nopagebreak[4]
\[
r=h\cot\beta.
\]

\gkanswer{答案：底面内切圆半径为 $h\cot\beta$}

\end{answer}
\begin{answer}
\gksolhead{16}{}

\textbf{题目}：

\Needspace{5\baselineskip}
已知圆

\nopagebreak[4]
\[
C:x^2+y^2+4x-12y+39=0
\]

\Needspace{5\baselineskip}
和直线

\nopagebreak[4]
\[
L:3x-4y+5=0.
\]

求圆 $C$ 关于直线 $L$ 的对称的圆的方程。

\Needspace{4\baselineskip}
\textbf{解析}

\Needspace{5\baselineskip}
圆方程配方为

\nopagebreak[4]
\[
(x+2)^2+(y-6)^2=1.
\]

圆心为 $O(-2,6)$，半径为 $1$。关于直线
\Needspace{5\baselineskip}
$3x-4y+5=0$ 的对称点坐标公式给出

\nopagebreak[4]
\[
\begin{aligned}
d&=3(-2)-4(6)+5=-25,\\
x'&=-2-\frac{2\cdot3d}{3^2+(-4)^2}=4,\\
y'&=6-\frac{2\cdot(-4)d}{3^2+(-4)^2}=-2.
\end{aligned}
\]

\Needspace{5\baselineskip}
故对称圆圆心为 $(4,-2)$、半径仍为 $1$，其方程为

\nopagebreak[4]
\[
(x-4)^2+(y+2)^2=1,
\]

\Needspace{5\baselineskip}
即

\nopagebreak[4]
\[
x^2+y^2-8x+4y+19=0.
\]

\gkanswer{答案：对称圆方程为 $(x-4)^2+(y+2)^2=1$，即 $x^2+y^2-8x+4y+19=0$}

\end{answer}
\begin{answer}
\gksolhead{17}{}

\textbf{题目}：

\Needspace{5\baselineskip}
设首项为 $1$、公比为 $q$（$q>0$）的等比数列的前 $n$ 项之和为 $S_n$，又设

\nopagebreak[4]
\[
T_n=\frac{S_n}{S_{n+1}},\qquad n=1,2,\ldots
\]

\Needspace{5\baselineskip}
求

\nopagebreak[4]
\[
\lim_{n\to\infty}T_n.
\]

\Needspace{4\baselineskip}
\textbf{解析}

\Needspace{5\baselineskip}
当 $q=1$ 时，$S_n=n$，所以

\nopagebreak[4]
\[
T_n=\frac n{n+1}\longrightarrow1.
\]

\Needspace{5\baselineskip}
当 $q\ne1$ 时，

\nopagebreak[4]
\[
S_n=\frac{q^n-1}{q-1},
\qquad
T_n=\frac{q^n-1}{q^{n+1}-1}.
\]

\Needspace{5\baselineskip}
若 $0<q<1$，令 $n\to\infty$ 得

\nopagebreak[4]
\[
T_n\longrightarrow\frac{-1}{-1}=1.
\]

\Needspace{5\baselineskip}
若 $q>1$，分子分母同除以 $q^{n+1}$，得

\nopagebreak[4]
\[
T_n=\frac{q^{-1}-q^{-(n+1)}}{1-q^{-(n+1)}}
\longrightarrow\frac1q.
\]

\Needspace{5\baselineskip}
因此

\nopagebreak[4]
\[
\lim_{n\to\infty}T_n=
\begin{cases}
1,&0<q\le1,\\[2pt]
\dfrac1q,&q>1.
\end{cases}
\]

\gkanswer{答案：$\displaystyle\lim_{n\to\infty}T_n=\begin{cases}1,&0<q\le1,\\[2pt]\dfrac1q,&q>1.\end{cases}$}

\end{answer}